\section*{Exercises}
\begin{exercise}\label{ex:5.1}Prove $\sum_{j=1}^{n}j=n(n+1)/2$ for $n\ge0$ by induction. Show the algebra in the transition.\end{exercise}
\begin{exercise}\label{ex:5.2}If $y_0=3$ and $y_{n+1}=2y_n$, find and prove a formula for $y_n$.\end{exercise}
\begin{exercise}\label{ex:5.3}Run Euclid's algorithm on $(1071,462)$. Record each quotient, remainder, and the final greatest common divisor.\end{exercise}
\begin{exercise}\label{ex:5.4}Starting with five plus signs and flipping two signs at a time, can you reach a configuration with exactly three minus signs? Justify without listing all moves.\end{exercise}
\begin{exercise}\label{ex:5.5}For the recurrence $x_{n+1}=(x_n+1)/2$ with arbitrary real $x_0$, prove the error formula $x_n-1=2^{-n}(x_0-1)$ by induction, writing out both the base case and the induction step. Explain why the error never becomes exactly zero when $x_0\ne1$.\end{exercise}
\begin{exercise}\label{ex:5.6}An induction proof uses $P(k)$ to prove $P(k+2)$ and checks only $P(0)$. Which indices has it actually covered? Give the additional starting case needed to cover all natural numbers.\end{exercise}
